\section{BÀI 4. XÁC ĐỊNH BẬC PHẢN ỨNG}
\subsection{THÍ NGHIỆM}
\subsubsection{Thí nghiệm 1: Xác định bậc phản ứng theo Na\textsubscript{2}S\textsubscript{2}O\textsubscript{3}}
\textbf{a. Tiến trình thí nghiệm}
\begin{itemize}
    \item Chuẩn bị 3 ống nghiệm chứa H\textsubscript{2}SO\textsubscript{4} và 3 bình tam giác chứa Na\textsubscript{2}S\textsubscript{2}O\textsubscript{3} và H\textsubscript{2}O theo bảng sau:
    \begin{table}[H]
        \centering
        \begin{tabular}{|c|c|c|c|}
            \hline
            \multirow{2}{*}{\textbf{TN}} & \textbf{Ống nghiệm} & \multicolumn{2}{c|}{\textbf{Erlen}} \\
            \cline{2-4}
            & \textbf{V(ml) H\textsubscript{2}SO\textsubscript{4} 0.4M} & \textbf{V(ml) Na\textsubscript{2}S\textsubscript{2}O\textsubscript{3} 0.1M} & \textbf{V(ml) H\textsubscript{2}O} \\
            \hline
            1 & 8 & 4 & 28 \\
            \hline
            2 & 8 & 8 & 24 \\
            \hline
            3 & 8 & 16 & 16 \\
            \hline
        \end{tabular}
    \end{table}
    \item Dùng Pipet vạch lấy axit cho vào ống nghiệm.
    \item Dùng Buret cho nước vào 3 Erlen.
    \item Tráng Buret bằng Na\textsubscript{2}S\textsubscript{2}O\textsubscript{3} 0.1M rồi tiếp tục dùng Buret cho Na\textsubscript{2}S\textsubscript{2}O\textsubscript{3} vào 3 Erlen.
    \item Lần lượt cho phản ứng từng cặp ống nghiệm và Erlen như sau:
    \begin{itemize}
        \item Đổ nhanh axit trong ống nghiệm vào Erlen.
        \item Bấm đồng hồ (khi 2 dung dịch tiếp xúc nhau).
        \item Lắc nhẹ, sau đó để yên quan sát, khi thấy dung dịch chuyển sang đục thì bấm dừng đồng hồ.
    \end{itemize}
    \item Lặp lại thí nghiệm lấy giá trị trung bình.
\end{itemize}

\textbf{b. Kết quả}
\begin{itemize}
    \item Bảng số liệu:
    \begin{table}[H]
        \centering
        \begin{tabular}{|c|c|c|c|c|c|}
            \hline
            \textbf{TN} & \textbf{[Na\textsubscript{2}S\textsubscript{2}O\textsubscript{3}]} & \textbf{[H\textsubscript{2}SO\textsubscript{4}]} & \textbf{$\Delta t_1$} & \textbf{$\Delta t_2$} & \textbf{$\Delta t_{TB}$} \\
            \hline
            1 & 0.01 & 0.08 & 142.1 & 130 & 136.05 \\
            \hline
            2 & 0.02 & 0.08 & 75.4 & 67.4 & 71.4 \\
            \hline
            3 & 0.04 & 0.08 & 33.2 & 45.6 & 39.4 \\
            \hline
        \end{tabular}
    \end{table}
    \item Ta có: $m = \frac{\log\left(\frac{\Delta t_1}{\Delta t_2}\right)}{\log(2)}$
    \item Từ TN1 và TN2 ta có: $m_1 = \frac{\log\left(\frac{\Delta t_1}{\Delta t_2}\right)}{\log(2)} = \frac{\log\left(\frac{136.05}{71.4}\right)}{\log(2)} \approx 0.93$
    \item Từ TN2 và TN3 ta có: $m_2 = \frac{\log\left(\frac{\Delta t_2}{\Delta t_3}\right)}{\log(2)} = \frac{\log\left(\frac{71.4}{39.4}\right)}{\log(2)} \approx 0.86$
    \item Bậc phản ứng theo Na\textsubscript{2}S\textsubscript{2}O\textsubscript{3}: $m = \frac{m_1 + m_2}{2} = \frac{0.93 + 0.86}{2} = 0.895$
\end{itemize}

\subsubsection{Thí nghiệm 2: Xác định bậc phản ứng theo H\textsubscript{2}SO\textsubscript{4}}
\textbf{a. Tiến trình thí nghiệm}
\begin{itemize}
    \item Thao tác tương tự phần 1 với lượng axit và Na\textsubscript{2}S\textsubscript{2}O\textsubscript{3} theo bảng sau:
    \begin{table}[H]
        \centering
        \begin{tabular}{|c|c|c|c|}
            \hline
            \multirow{2}{*}{\textbf{TN}} & \textbf{Ống nghiệm} & \multicolumn{2}{c|}{\textbf{Erlen}} \\
            \cline{2-4}
            & \textbf{V(ml) H\textsubscript{2}SO\textsubscript{4} 0.4M} & \textbf{V(ml) Na\textsubscript{2}S\textsubscript{2}O\textsubscript{3} 0.1M} & \textbf{V(ml) H\textsubscript{2}O} \\
            \hline
            1 & 4 & 8 & 1 \\
            \hline
            2 & 8 & 8 & 2 \\
            \hline
            3 & 16 & 8 & 3 \\
            \hline
        \end{tabular}
    \end{table}
\end{itemize}

\textbf{b. Kết quả}
\begin{itemize}
    \item Bảng số liệu:
    \begin{table}[H]
        \centering
        \begin{tabular}{|c|c|c|c|c|c|}
            \hline
            \textbf{TN} & \textbf{[Na\textsubscript{2}S\textsubscript{2}O\textsubscript{3}]} & \textbf{[H\textsubscript{2}SO\textsubscript{4}]} & \textbf{$\Delta t_1$} & \textbf{$\Delta t_2$} & \textbf{$\Delta t_{TB}$} \\
            \hline
            1 & 0.02 & 0.04 & 73.6 & 73.4 & 73.5 \\
            \hline
            2 & 0.02 & 0.08 & 70 & 65.9 & 67.95 \\
            \hline
            3 & 0.02 & 0.16 & 56.3 & 58.5 & 57.4 \\
            \hline
        \end{tabular}
    \end{table}
    \item Từ TN1 và TN2 ta có: $n_1 = \frac{\log\left(\frac{\Delta t_1}{\Delta t_2}\right)}{\log(2)} = \frac{\log\left(\frac{73.5}{67.95}\right)}{\log(2)} \approx 0.11$
    \item Từ TN2 và TN3 ta có: $n_2 = \frac{\log\left(\frac{\Delta t_2}{\Delta t_3}\right)}{\log(2)} = \frac{\log\left(\frac{67.95}{57.4}\right)}{\log(2)} \approx 0.24$
    \item Bậc phản ứng theo H\textsubscript{2}SO\textsubscript{4}: $n = \frac{n_1 + n_2}{2} = \frac{0.11 + 0.24}{2} = 0.175$
\end{itemize}

\subsection{TRẢ LỜI CÂU HỎI}
\begin{enumerate}
    \item \textbf{Trong thí nghiệm trên nồng độ [Na\textsubscript{2}S\textsubscript{2}O\textsubscript{3}] và của [H\textsubscript{2}SO\textsubscript{4}] đã ảnh hưởng thế nào lên vận tốc phản ứng. Viết lại biểu thức tính vận tốc phản ứng. Xác định bậc của phản ứng?}
    \begin{itemize}
        \item Nồng độ [Na\textsubscript{2}S\textsubscript{2}O\textsubscript{3}] tỉ lệ thuận với tốc độ phản ứng.
        \item Nồng độ [H\textsubscript{2}SO\textsubscript{4}] hầu như không ảnh hưởng tới tốc độ phản ứng.
        \item Biểu thức tính vận tốc: $V = k \cdot [\text{Na}_2\text{S}_2\text{O}_3]^{0.9} \cdot [\text{H}_2\text{SO}_4]^{0.07}$
        \item Bậc của phản ứng: $0.9 + 0.07 = 0.97$
    \end{itemize}
    \item \textbf{Cơ chế của phản ứng trên có thể được viết như sau:}
    \begin{align}
        &\text{H}_2\text{SO}_4 + \text{Na}_2\text{S}_2\text{O}_3 \rightarrow \text{Na}_2\text{SO}_4 + \text{H}_2\text{S}_2\text{O}_3 \quad &(1) \nonumber \\
        &\text{H}_2\text{S}_2\text{O}_3 \rightarrow \text{H}_2\text{SO}_3 + \text{S}\downarrow \quad &(2) \nonumber
    \end{align}
    \textbf{Dựa vào kết quả thí nghiệm có thể kết luận phản ứng (1) hay (2) là phản ứng quyết định vận tốc phản ứng tức là phản ứng xảy ra nhanh hay chậm không? Tại sao? Lưu ý trong các phản ứng trên, lượng axit [H\textsubscript{2}SO\textsubscript{4}] luôn dư so với [Na\textsubscript{2}S\textsubscript{2}O\textsubscript{3}].}
    \begin{itemize}
        \item (1) Là phản ứng trao đổi ion nên tốc độ phản ứng xảy ra rất nhanh.
        \item (2) Là phản ứng tự oxi khử nên tốc độ phản ứng xảy ra rất chậm.
    \end{itemize}
    \item \textbf{Dựa trên cơ sở của phương pháp thí nghiệm thì vận tốc xác định được trong các thí nghiệm được xem là vận tốc trung bình hay vận tốc tức thời?}
    Vận tốc xác định bằng $\frac{\Delta C}{\Delta t}$ vì $\Delta C \approx 0$ (biến thiên nồng độ của lưu huỳnh không đáng kể trong khoảng thời gian $\Delta t$) nên vận tốc trong các thí nghiệm trên được xem là vận tốc tức thời.
    \item \textbf{Thay đổi thứ tự cho [H\textsubscript{2}SO\textsubscript{4}] và [Na\textsubscript{2}S\textsubscript{2}O\textsubscript{3}] thì bậc phản ứng có thay đổi không. Tại sao?}
    Thay đổi thứ tự cho [H\textsubscript{2}SO\textsubscript{4}] và [Na\textsubscript{2}S\textsubscript{2}O\textsubscript{3}] thì bậc phản ứng không thay đổi. Vì ở một nhiệt độ xác định bậc phản ứng chỉ phụ thuộc vào bản chất của hệ (nồng độ, nhiệt độ, diện tích bề mặt, áp suất) mà không phụ thuộc vào thứ tự chất phản ứng.
\end{enumerate}
